\documentclass[11pt,a4paper]{article}
\usepackage[utf8]{inputenc}
\usepackage[spanish,es-noquoting]{babel}
\usepackage{amsmath,amsfonts,amssymb}
\usepackage{geometry}
\geometry{top=2.5cm, bottom=2.5cm, left=2.5cm, right=2.5cm}
\usepackage{hyperref}
\usepackage{booktabs}
\usepackage{tabularx}
\usepackage{xcolor}
\usepackage{tcolorbox}
\usepackage{enumitem}
\usepackage{graphicx}

% Colores institucionales y de diseño
\definecolor{primary}{RGB}{15, 32, 67}
\definecolor{secondary}{RGB}{0, 114, 206}
\definecolor{accentgreen}{RGB}{46, 125, 50}
\definecolor{lightgray}{RGB}{245, 247, 250}

\hypersetup{
    colorlinks=true,
    linkcolor=secondary,
    urlcolor=secondary,
    pdftitle={Portafolio de Evidencias No. 4 - Wheels Externado}
}

\begin{document}

% ---------------------------------------------------------
% PORTADA INSTITUCIONAL
% ---------------------------------------------------------
\begin{titlepage}
    \centering
    \vspace*{1cm}
    {\Huge \textbf{UNIVERSIDAD EXTERNADO DE COLOMBIA}}\\[0.4cm]
    {\Large \textbf{PREGRADO EN CIENCIA DE DATOS --- EMPRENDIMIENTO}}\\[1.8cm]
    
    \rule{\linewidth}{0.8mm}\\[0.4cm]
    {\Huge \textbf{PORTAFOLIO DE EVIDENCIAS NO. 4}}\\[0.3cm]
    {\Large \textbf{Capturar Valor: Estructura de Costos, Fuentes de Ingreso y Unit Economics}}\\[0.4cm]
    \rule{\linewidth}{0.8mm}\\[1.8cm]
    
    \begin{tcolorbox}[colback=primary!5,colframe=primary]
        \centering \Huge
        \textbf{Wheels Externado} \\[0.2cm]
        \large \textit{Plataforma de Movilidad y Carpooling Universitario Inteligente}
    \end{tcolorbox}
    
    \vspace{2cm}
    
    \textbf{Grupo 5:}\\[0.3cm]
    \begin{tabular}{ll}
        \textbf{Santiago Sandoval} & \textit{Product \& Algorithm Lead} \\
        \textbf{Sebastián Ramos} & \textit{Director de Estrategia \& Modelo Financiero} \\
        \textbf{Julián Duarte} & \textit{Arquitectura de Datos \& Backend} \\
        \textbf{Tomás Rincón} & \textit{Operaciones, Seguridad \& Compliance} \\
        \textbf{Julián Jiménez} & \textit{UX/UI \& Mobile App Engineering} \\
    \end{tabular}
    
    \vfill
    {\large Semestre 2026-II}\\[0.2cm]
    {\normalsize Bogot\'a D.C., Colombia}
\end{titlepage}

\newpage
\tableofcontents
\newpage

% ---------------------------------------------------------
% SECCIÓN 1: NOTA METODOLÓGICA DEL PIVOTE ESTRATÉGICO
% ---------------------------------------------------------
\section{Declaraci\'on y Fundamentaci\'on del Pivote Estrat\'egico}

Durante las primeras tres fases del curso (Portafolios 1 al 3 y Entrega 1), el equipo formul\'o y valid\'o la propuesta de \textit{StatCredit AI}, un motor anal\'itico B2B de scoring crediticio alternativo para microempresarios e independientes. A pesar de su solidez conceptual y rigor regulatorio (normativa SARC), el equipo identific\'o fricciones cr\'iticas de validaci\'on pr\'actica: ciclos de venta corporativos extremadamente largos en el sector bancario (9 a 18 meses) y barreras de acceso a bases de datos transaccionales reales.

Siguiendo las ense\~nanzas de Blank (2021) y Bianchi \& Verganti (2021), un equipo emprendedor debe tener la disciplina de pivotear cuando identifica un dolor m\'as urgente, cotidiano, medible y accesible en su entorno inmediato.

\begin{tcolorbox}[colback=secondary!5,colframe=secondary,title=\textbf{Nuevo Enunciado del Proyecto: Wheels Externado}]
Desarrollo de una plataforma digital y aplicaci\'on m\'ovil de viajes compartidos (\textit{carpooling} universitario) dise\~nada exclusivamente para la comunidad de la \textbf{Universidad Externado de Colombia}, orientada a sustituir la actual red informal e insegura de grupos de WhatsApp por un ecosistema con verificaci\'on institucional obligatoria, algoritmos de coincidencia geoespacial de rutas, pagos automatizados y garant\'ias de seguridad en ruta.
\end{tcolorbox}

% ---------------------------------------------------------
% SECCIÓN 2: EL PROBLEMA Y LA OPORTUNIDAD EN EL EXTERNADO
% ---------------------------------------------------------
\section{Caracterizaci\'on del Problema: La Realidad de los Grupos de WhatsApp}

La comunidad de la Universidad Externado de Colombia (aproximadamente 12.000 estudiantes de pregrado y posgrado, m\'as de 1.500 docentes y personal administrativo) enfrenta un desaf\'io estructural de movilidad:
\begin{enumerate}[leftmargin=*]
    \item \textbf{Topograf\'ia y Localizaci\'on:} El campus est\'a enclavado en los Cerros Orientales (barrio Egipto / La Candelaria). La subida a pie desde las estaciones de TransMilenio (Las Aguas y Universidades) representa una pendiente pronunciada y es percibida como insegura en jornadas nocturnas.
    \item \textbf{El Fen\'omeno Informal de los ``Wheels'':} Para mitigar los costos de transporte y el tiempo de viaje, han surgido espont\'aneamente m\'as de 15 grupos masivos de WhatsApp segmentados por zonas geogr\'aficas (Ch\'ia, Norte, Calle 80, Suba, Teusaquillo, Sur/Soacha).
    \item \textbf{Inseguridad Cr\'itica:} Cualquier persona que obtenga el enlace de invitaci\'on ingresa a los grupos. No hay verificaci\'on de identidad institucional, ni constancia del estado legal de los veh\'iculos (SOAT, revisi\'on t\'ecnico-mec\'anica o antecedentes).
    \item \textbf{Caos Operativo y Saturaci\'on:} En horas pico, los mensajes de oferta y demanda de cupos se pierden en segundos, generando cancelaciones imprevistas, descoordinaci\'on en puntos de recogida y fricciones por pagos en efectivo o transferencias manuales.
\end{enumerate}

% ---------------------------------------------------------
% SECCIÓN 3: FUENTES DE INGRESO Y MODELO DE MONETIZACIÓN
% ---------------------------------------------------------
\section{Modelo de Captura de Valor (Fuentes de Ingreso --- Aulet Cap. 15)}

A diferencia de las aplicaciones de transporte comercial individual, \textit{Wheels Externado} se fundamenta en un modelo de \textbf{movilidad colaborativa institucional}, donde los usuarios comparten los costos de viaje (combustible, peajes, parqueadero). La plataforma captura valor a trav\'es de tres mecanismos complementarios:

\begin{table}[ht!]
\centering
\small
\begin{tabularx}{\textwidth}{p{4.2cm} X p{3.8cm}}
\toprule
\textbf{L\'inea de Ingreso} & \textbf{Mecanismo de Cobro} & \textbf{Tarifa Estimada (COP)} \\
\midrule
\textbf{1. Cuota Operativa por Reserva} & Tarifa fija digital cobrada por cupo completado. & \$600 -- \$800 COP por viaje. \\
\addlinespace
\textbf{2. Membres\'ia Conductor Frecuente} & Suscripci\'on mensual sin comisiones por viaje y visibilidad destacada. & \$14.900 COP / mes. \\
\addlinespace
\textbf{3. Convenio Institucional} & Tarifa de soporte y anal\'itica de movilidad y huella de carbono. & Pago semestral institucional. \\
\bottomrule
\end{tabularx}
\caption{Estructura de fuentes de ingreso de Wheels Externado.}
\end{table}

% ---------------------------------------------------------
% SECCIÓN 4: ESTRUCTURA DE COSTOS Y UNIT ECONOMICS
% ---------------------------------------------------------
\section{Estructura de Costos y Unit Economics (Aulet Pasos 17 y 19)}

\subsection{Costos Fijos Operativos (Mensuales)}
\begin{itemize}[leftmargin=*]
    \item \textbf{Infraestructura en la Nube (AWS / Google Cloud):} Servidores de aplicaci\'on, bases de datos PostgreSQL con PostGIS para consultas espaciales: \$180 USD/mes (\$720.000 COP/mes).
    \item \textbf{APIs de Mapas y Ruteo (Mapbox / OpenStreetMap):} Optimizaci\'on de costos utilizando capas abiertas con cach\'e local para c\'alculo de matrices de distancia: \$120 USD/mes (\$480.000 COP/mes).
    \item \textbf{Mantenimiento y Soporte T\'ecnico:} Monitoreo continuo y resoluci\'on de incidencias: \$1.200.000 COP/mes.
\end{itemize}

\subsection{Costos Variables por Transacci\'on}
\begin{itemize}[leftmargin=*]
    \item \textbf{Pasarela de Pagos (Wompi / PayU):} Costo porcentual + fijo por dispersi\'on hacia cuentas Nequi/Daviplata (\$200 COP + 1,5\% por transacci\'on).
    \item \textbf{Validaci\'on de Datos en RUNT (Onboarding):} Consulta de antecedentes y SOAT por conductor registrado (\$1.500 COP costo \'unico por verificaci\'on).
\end{itemize}

\subsection{M\'etricas Unitarias Clave (Unit Economics)}
\begin{tcolorbox}[colback=accentgreen!5,colframe=accentgreen,title=\textbf{Indicadores de Salud Financiera Unitaria}]
\begin{itemize}
    \item \textbf{Costo de Adquisici\'on de Clientes (CAC):} \$3.800 COP por usuario activo (apalancado en marketing org\'anico en campus, activaciones en facultades y migraci\'on de los grupos de WhatsApp existentes).
    \item \textbf{Valor del Tiempo de Vida del Cliente (LTV):} Con un promedio de 8 viajes semanales (32 viajes/mes) a una comisi\'on promedio de \$700 COP por viaje y una retenci\'on m\'inima de 8 meses por a\~no acad\'emico:
    \[ \text{LTV} = 32 \times 700 \times 8 \times (1 - 0.25) \approx \$134.400 \text{ COP} \]
    \item \textbf{Relaci\'on LTV / CAC:}
    \[ \frac{\text{LTV}}{\text{CAC}} = \frac{134.400}{3.800} \approx 35.3 \times \]
    Supera con amplia holgura la recomendaci\'on de Aulet ($\text{LTV} \ge 3\times \text{CAC}$), gracias al costo marginal pr\'acticamente nulo de adquisici\'on dentro de un campus universitario cerrado.
\end{itemize}
\end{tcolorbox}

% ---------------------------------------------------------
% SECCIÓN 5: LIENZO CANVAS DE VALOR
% ---------------------------------------------------------
\section{Lienzo Canvas Estrat\'egico: Zona Capturar Valor}

\begin{tcolorbox}[colback=primary!5,colframe=primary,title=\textbf{S\'intesis de la Zona Capturar Valor}]
\begin{description}
    \item[\textbf{Estructura de Costos:}] Mantenimiento de servidores en la nube, licencias de geolocalizaci\'on, comisiones de dispersi\'on bancaria, campa\~nas presenciales de activaci\'on en el campus y validaci\'on inicial de veh\'iculos.
    \item[\textbf{Fuentes de Ingresos:}] Comisi\'on por viaje coordinado y completado (\$600--\$800 COP), planes premium para conductores con beneficios de parqueadero, y alianzas con comercios del centro hist\'orico para promociones cruzadas.
\end{description}
\end{tcolorbox}

\vspace{0.5cm}
\begin{center}
\small \textit{Universidad Externado de Colombia --- Pregrado en Ciencia de Datos --- Curso de Emprendimiento 2026}
\end{center}

\end{document}
